\begin{thebibliography}{99}
\addcontentsline{toc}{section}{References}

\bibitem{persiflage}
Frank Calegari (Persiflage),
\emph{\(\zeta(5)\) is irrational}, 24 September 2026.
\href{https://galoisrepresentations.org/2026/09/24/zeta5-is-irrational/}{Blog post}.
The guiding exposition for this draft. The displayed post date is earlier
than Lai's write-up; no dependence of the post on Lai is asserted.
Calegari identifies Persiflage as his blog on his
\href{https://math.uchicago.edu/~fcale/research.html}{research page}.

\bibitem{fauzan}
Aabir Fauzan,
\emph{\(\zeta(5)\) is irrational}, manuscript dated 17 September 2026,
32 pages, Zenodo record 22826419.
\href{https://zenodo.org/records/22826419}{DOI: 10.5281/zenodo.22826419}.
The determinant, factorial normalization, local arithmetic strategy,
comparison measure, and analytic certificate originate here.

\bibitem{firsching}
Moritz Firsching and contributors,
\emph{Zeta5: Lean formalization of the irrationality of \(\zeta(5)\)}.
\href{https://github.com/mo271/Zeta5/tree/7fe736760f4b96bfdb4334b68e3b3124ecbe10b0}{Repository at commit
\texttt{7fe7367}}, 27 September 2026. Apache License 2.0.
The substantive initial formalization is commit
\href{https://github.com/mo271/Zeta5/commit/2c80b320d06b4395605edbd413649d19094507af}{\texttt{2c80b32}}.
The revised local estimates, cutoff 400, full rank loss, and global
normalization used in this draft are translated from the
\texttt{Apery/Arith}, \texttt{Apery/Norm.lean}, and
\texttt{Apery/Growth} files, with exact locations in the source ledger.
The formal development depends on Lean, Mathlib and
\href{https://github.com/mo271/PrimeNumberTheoremAnd/tree/04d81b2a68ea3548b9990813fdc463c78511f412}{PrimeNumberTheoremAnd};
their contributors retain credit for those libraries and formalized
background theorems. Paul-Lez is recorded as a coauthor of the Comparator
CI addition. No local rebuild of the complete Lean development is claimed.

\bibitem{lai}
Chun Wing Lai,
\emph{On Fauzan's Proof that \(\zeta(5)\) is Irrational},
v1, 29 September 2026.
\href{https://ai.vixra.org/pdf/2609.0080v1.pdf}{ai.vixra:2609.0080v1}.
Used as a later exposition of the construction and as a cross-check;
its printed arithmetic conclusion is not an input to this draft.

\bibitem{prevost}
Marc Pr\'evost,
\emph{Remainder Pad\'e approximants for the Hurwitz zeta function},
Results in Mathematics \textbf{74} (2019), article 51.
\href{https://arxiv.org/pdf/1709.05389v1}{arXiv:1709.05389v1},
15 September 2017. Theorems 1--2, equations (3.1)--(3.4), and Section 4.1
give the positive representation and moments used here.

\bibitem{prevost1996}
Marc Pr\'evost,
\emph{A new proof of the irrationality of \(\zeta(2)\) and \(\zeta(3)\)
using Pad\'e approximants},
Journal of Computational and Applied Mathematics \textbf{67} (1996),
219--235.
\href{https://doi.org/10.1016/0377-0427(95)00019-4}{DOI: 10.1016/0377-0427(95)00019-4}.
Historical context; the present proof does not invoke an unproved
result from this paper.

\bibitem{prevost-rivoal}
Marc Pr\'evost and Tanguy Rivoal,
\emph{Diagonal convergence of the remainder Pad\'e approximants for the
Hurwitz zeta function},
Journal of Number Theory \textbf{222} (2021), 346--361.
\href{https://doi.org/10.1016/j.jnt.2020.10.019}{DOI: 10.1016/j.jnt.2020.10.019}.
Background on approximation; no quantitative bound from this paper is used.

\bibitem{zudilin}
Wadim Zudilin,
\emph{A determinantal approach to irrationality},
Constructive Approximation \textbf{45} (2017), 301--310.
\href{https://arxiv.org/pdf/1507.05697v2}{arXiv:1507.05697v2},
especially Proposition 2 and Section 2.

\bibitem{uvarov}
V.\,B. Uvarov,
\emph{The connection between systems of polynomials that are orthogonal
with respect to different distribution functions},
USSR Computational Mathematics and Mathematical Physics
\textbf{9}, no.\ 6 (1969), 25--36.
\href{https://doi.org/10.1016/0041-5553(69)90124-4}{DOI: 10.1016/0041-5553(69)90124-4}.
Historical attribution through Krattenthaler's discussion \cite{krattenthaler};
the original article was not used as a proof source.

\bibitem{krattenthaler}
Christian Krattenthaler,
\emph{A determinant identity for moments of orthogonal polynomials that
implies Uvarov's formula for the orthogonal polynomials of rationally
related densities},
\href{https://arxiv.org/pdf/2103.03969v1}{arXiv:2103.03969v1}, 2021.
Introduction and Theorem 1 provide the general setting for rational
modification of a measure. The theorem itself is not required below.

\bibitem{brown}
Francis Brown,
\emph{Mellin transforms, transfinite diameter and rational approximations
of integrals},
\href{https://arxiv.org/html/2604.20741v1}{arXiv:2604.20741v1}, 2026.
Sections 3.1--3.2 and 7.4--7.5 give related pairing and normalization
frameworks. This is contextual credit, not a claim that Brown states
the specific estimates proved here.

\bibitem{andreief}
C. Andr\'eief,
\emph{Note sur une relation entre les int\'egrales d\'efinies des produits
des fonctions}, M\'emoires de la Soci\'et\'e des Sciences physiques et
naturelles de Bordeaux, third series, \textbf{2} (1886), 1--14.
For the identity and publication history we consulted Peter J. Forrester,
\emph{Meet Andr\'eief, Bordeaux 1886, and Andreev, Kharkov 1882--83},
\href{https://arxiv.org/pdf/1806.10411v1}{arXiv:1806.10411v1}, 2018.
The identity is proved in the text.

\bibitem{dlmf}
NIST Digital Library of Mathematical Functions,
\href{https://dlmf.nist.gov/4.36.E3}{4.36.3} (partial fractions of \(\coth\)),
\href{https://dlmf.nist.gov/25.6.E2}{25.6.2} (Euler's even-zeta formula),
\href{https://dlmf.nist.gov/25.11.E29}{25.11.29} (Hermite's formula),
\href{https://dlmf.nist.gov/24.2}{24.2} (Bernoulli generating functions),
\href{https://dlmf.nist.gov/24.4}{24.4} (Bernoulli identities), and
\href{https://dlmf.nist.gov/24.10.i}{24.10(i)}
(von Staudt--Clausen). Consulted 14 November 2026.

\end{thebibliography}
